\begin{table}[tbp]
\caption{Transition actions and the continuity claims they can support.}
\label{tab:ch8-transition}
\centering\small\renewcommand{\arraystretch}{1.13}
\begin{tabularx}{\linewidth}{@{}>{\raggedright\arraybackslash}p{29mm}>{\raggedright\arraybackslash}X>{\raggedright\arraybackslash}X@{}}\toprule
Observed action & Immediate evidential contribution & Further evidence for continuity\\\midrule
Recommend an alternative & Identifies a possible destination. & Fit to required functions and compatibility with affected users' environments.\\
Publish a migration guide & Describes a transition procedure. & Successful use of the procedure, effort, failures, and residual dependence.\\
Link a fork & Identifies a continuing branch and potential maintainers. & Independent maintenance, accepted responsibility, releases, and response capacity.\\
Resolve or close requests & Documents a disposition of outstanding work. & Whether defects were remedied, transferred, declined with explanation, or left unresolved.\\
Set a timetable & Makes stages and deadlines explicit. & Whether users received notice and could complete the required transition.\\
Provide contact details & Preserves a possible communication route. & Availability, agreed scope, and actual handling of subsequent requests.\\\bottomrule
\end{tabularx}
\tablenote{Book synthesis from the practices in Table~\ref{tab:ch8-practices}. The final column states evaluation requirements; it does not report outcomes measured in the archive study.}
\end{table}
